\begin{proof}
\begin{enumerate}
\item \textbf{Degenerate cases.}  
If \(n=1\) the statement reads \(a_{1}b_{1}\ge a_{1}b_{1}\), which holds with equality.  
If the \(a\)-sequence is constant, say \(a_{i}=c\) for all \(i\), then
\[
n\sum_{i=1}^{n}a_i b_i
   = n c\sum_{i=1}^{n} b_i
   =\Bigl(\sum_{i=1}^{n} a_i\Bigr)\Bigl(\sum_{i=1}^{n} b_i\Bigr),
\]
so the inequality is an equality.  The same argument applies when the \(b\)-sequence is constant.  
Hence we may assume \(n\ge2\) and that neither sequence is identically constant.

\item \textbf{Pairwise non‑negativity.}  
For any indices \(i<j\) we have \(a_i\le a_j\) and \(b_i\le b_j\); therefore
\[
(a_i-a_j)(b_i-b_j)\ge0. \tag{1}
\]

\item \textbf{Summing over all pairs.}  
Summing \((1)\) over all unordered pairs \((i,j)\) with \(1\le i<j\le n\) gives
\[
\sum_{1\le i<j\le n}(a_i-a_j)(b_i-b_j)\ge0. \tag{2}
\]

\item \textbf{Algebraic identity.}  
For a fixed pair \((i,j)\),
\[
(a_i-a_j)(b_i-b_j)=a_i b_i + a_j b_j - a_i b_j - a_j b_i .
\]
Summing this expression over all \(i<j\) yields
\begin{align*}
\sum_{i<j}(a_i-a_j)(b_i-b_j)
&= \underbrace{\sum_{i<j}(a_i b_i + a_j b_j)}_{\text{each }a_k b_k\text{ appears }n-1\text{ times}}
   -\underbrace{\sum_{i<j}(a_i b_j + a_j b_i)}_{\text{each }a_k b_\ell\;(k\neq\ell)\text{ appears once}} \\[4pt]
&= (n-1)\sum_{k=1}^{n}a_k b_k
   -\Bigl(\sum_{k=1}^{n}a_k\Bigr)\Bigl(\sum_{\ell=1}^{n}b_\ell\Bigr)
   +\sum_{k=1}^{n}a_k b_k \\[4pt]
&= n\sum_{k=1}^{n}a_k b_k
   -\Bigl(\sum_{k=1}^{n}a_k\Bigr)\Bigl(\sum_{\ell=1}^{n}b_\ell\Bigr). \tag{3}
\end{align*}
The counting argument above is valid for every integer \(n\ge2\).

\item \textbf{Chebyshev’s inequality.}  
Insert \((3)\) into \((2)\):
\[
0\le
n\sum_{i=1}^{n}a_i b_i
-\Bigl(\sum_{i=1}^{n}a_i\Bigr)\Bigl(\sum_{i=1}^{n}b_i\Bigr).
\]
Rearranging gives the desired inequality
\[
n\sum_{i=1}^{n}a_i b_i\;\ge\;
\Bigl(\sum_{i=1}^{n}a_i\Bigr)\Bigl(\sum_{i=1}^{n}b_i\Bigr). \tag{4}
\]

\item \textbf{Equality condition.}  
Equality in \((4)\) can occur only if equality holds in every instance of \((1)\), because \((2)\) is a sum of non‑negative terms.  
For a particular pair \((i,j)\) we have equality in \((1)\) precisely when \(a_i=a_j\) \emph{or} \(b_i=b_j\).  
Thus equality in the whole sum \((2)\) forces either

\begin{itemize}
\item \(a_i=a_j\) for all \(i,j\) (the \(a\)-sequence is constant), or
\item \(b_i=b_j\) for all \(i,j\) (the \(b\)-sequence is constant).
\end{itemize}
The case \(n=1\) is included in both possibilities. Hence equality holds \emph{iff} one of the two sequences is constant.

\item \textbf{Conclusion.}  
For any integer \(n\ge1\) and any real numbers
\[
a_{1}\le a_{2}\le\cdots\le a_{n},\qquad
b_{1}\le b_{2}\le\cdots\le b_{n},
\]
the inequality
\[
n\sum_{i=1}^{n} a_i b_i\;\ge\;
\Bigl(\sum_{i=1}^{n} a_i\Bigr)\Bigl(\sum_{i=1}^{n} b_i\Bigr)
\]
holds, with equality exactly when the \(a\)-sequence or the \(b\)-sequence is constant (including the trivial case \(n=1\)). \(\square\)
\end{enumerate}
\end{proof}